% Generated by src/tools/captions.py from evidence/figures/<fig>.csv and <fig>-caption.txt. Do not edit.
\newcommand{\CaptionPapyrus}{Papyrus of PHerc.~0826, seen in one flat cut across the scroll's axis, 6 mm on a side, around the largest fully covered square of the whole chain. (a) The raw scan, with the 8 of the 10 sheets grown from seed 2604 that cross this part of the cut. Sheet 0 holds that square, 17.7995 mm on a side, the largest over the 1,580 delivered sheets of 158 seeds; yellow is where the square crosses the cut. The blue lines are the seed's other sheets. The sheets follow the bright layers of papyrus. (b) The same cut with the surface prediction the growth runs on shown in green: it marks almost every layer, and each sheet in (a) is one path chosen through it. The scale bars are 1 mm. Only the scan and the shape of the sheets are shown; nothing is read along a sheet. Drawn by \texttt{src/tools/c-f10-papyrus.py}.}
\newcommand{\CaptionAtwoFlags}{One delivered sheet of PHerc.~0826 laid out in its own grid, as the growth built it: seed 5728, sheet 0. Grey is the part of the grid the sheet covers. An orange point is a place, checked every few cells, where the sheet meets the surface prediction at a steep angle, more than 25 degrees, which is how a sheet that jumps from one layer of papyrus to the next shows itself. Where such points gather into a patch of more than 101 points, they turn red, and the rule cuts the cells right around them (dark red, visible in the enlargement b). The blue outline is the largest fully covered square of the sheet as delivered, 11.7496 mm on a side; the green outline is the largest square left once the cut cells are taken out, 8.1948 mm. Of the 1,580 delivered sheets checked this way, this is the one where the rule costs the most millimetres, and it lowers the square on 20 of them. The scale bars are 5 mm. Only the sheet's shape is drawn; nothing is read along it. Drawn by \texttt{src/tools/c-f12-a2-flags.py}.}
\newcommand{\CaptionSetSeed}{Two runs from the same starting point of PHerc.~0826 (seed 6365), on the same cut across the scroll, the same 20 mm area and the same scale. The yellow ring is the starting point. (a) The run as the chain delivered it: its 10 sheets, drawn where they pass through this cut; 10 of them do, spreading over many layers of papyrus around the start. The red one holds the run's largest square, 10.2452 mm. (b) The same growth with one change, SetSeed, which starts the seed cells of each new patch at rest; the rest of the pipeline is unchanged. Of its 10 sheets, 7 pass through this cut, and they stay close to the layers that run through the starting point. Its largest square, 8.0756 mm, lies on a sheet that does not reach this cut, so no red line appears in (b). The scale bars are 5 mm. Only the scan and the shape of the sheets are shown; nothing is read along a sheet. Drawn by \texttt{src/tools/c-f13-setseed.py}.}
\newcommand{\CaptionCertifiedSquare}{The largest certified square on the chain's own sheets: the certified square of PHerc.~0826, on sheet 0 of seed 5364, drawn on the raw scan texture sampled along the sheet: raw scan, no ink detector. The sky blue outline is the certified square, 15.3693 mm on a side (411 cells), every cell of it covered and crossing none of the sheets it is compared with (the area study's sheets and the other sheets of its seed). No cell of it is in a self-conflict or in a cluster the steep angle rule cuts. It is the largest certified square of any of the 158 seeds of the chain as run measured for it, the next seed reaching 12.8975 mm. The orange outline is the hole free square of the same sheet, 17.1642 mm. At right, the 1.98 mm window framed in black, the one of highest local contrast inside the certified square, resampled from the raw scan at 8 points per cell side. The scale bars are 5 and 0.5 mm. Drawn by \texttt{src/tools/c-f16-certified-square.py}.}
\newcommand{\CaptionWhereSquares}{Where the best squares of PHerc.~0826 sit relative to each other, on the raw scan with no ink detector: the part of one axial cut around them, through the certified square of Figure~\ref{fig:certified}, read at level 1. The green lines are where the delivered sheets of the marked seeds cross the cut. Each mark is the centre of one square, labelled with its side. Sky blue marks the 3 largest certified squares and open orange the 5 largest hole free squares of the chain as run. The orange diamond is the largest square of the longer growths of the capped seeds, which is hole free and not certified one lamina. Each centre is drawn at its own position, projected along the height onto this cut. Drawn by \texttt{src/tools/c-f17-where-squares.py}.}
\newcommand{\CaptionBestWindows}{The three best clean windows of PHerc.~0826 on the raw scan, no ink detector, each at least 20 mm on a side on one delivered sheet and chosen by the checks of Section~\ref{sec:checks} with no person in the loop. They are seed 5364 sheet 1, seed 6473 sheet 0 and seed 6731 sheet 0. Their covered shares are 0.956311, 0.977619 and 0.987150, and their flagged shares 0, 0.000576 and 0.001712. Their fibre scores are 0.0843, 0.0799 and 0.0999. (a) to (c) The raw scan at each grown point, one pixel per cell of the sheet's own grid; the grid axis nearer the scan's height axis lies 14.4 to 30.8 degrees from it, so the fibres run slanted. (d) to (f) The same surfaces resampled so that rows follow constant height in the scan and columns follow equal arc length along the sheet, by the window search's own resampling tool, align.py, not by the renderer of the organisers used for the ink reads. (g) The straightened section of the first window along its j line. Each column is the raw scan along one grown point's normal, 60 voxels on either side, with the sheet on the middle row (faint sky blue line), and the rows are drawn 4 times their true scale. Within 15 voxels of the sheet, the brightest voxel lies a median of 6 voxels from it and within 3 voxels in a share 0.3701 of the columns, against 8 voxels and 0.2258 for a peak placed at random. That separates it from chance, but it does not prove one layer. The sections of the other two windows do not separate (medians of 8 to 10 voxels) and are not drawn. Scale bars: 5 mm, and 0.5 mm across the section. Drawn by \texttt{src/tools/c-f20-best-windows.py}.}
\newcommand{\CaptionRtwoC}{The certified square of the headline surface on PHerc.~0826, 28.1285 mm, with crossings between traced surfaces adjudicated, on the raw scan, no ink detector. Our search chose the start, the centre of the certified square of seed 6273; the organisers' tracer grew the surface; our checks ran on the whole crop of Section~\ref{sec:checks}. (a) The raw scan over the square at one pixel per node of the tracer's own grid, not turned: the scan's height already points up. That grid lies about 19.65 degrees in plane from the nearest axis aligned grid, so the fibres run slanted. Panels (b) and (c) are of the square that passed the checks against our sheets before the adjudication, 29.8961 mm. (b) An axial cut through that square's centre, with the tracer surface inside it and our sheet seed1045-S0 (grown by the chain); a share 0.7003 of that square's nodes lie within 3 voxels of that sheet. (c) The straightened sections along its two centre lines: the brightest layer lies within 3 voxels of the surface in a share 0.4553 of the columns along i and 0.2692 along j. The panels follow the route study's own figure tool, redrawn in this work. Drawn by \texttt{src/tools/c-f21-r2c.py}.}
\newcommand{\CaptionSquareFull}{The certified square of the headline surface on PHerc.~0826 at full resolution: the raw scan sampled along the surface at one sample per voxel (\FullSquareVoxel{}~$\mu$m), \FullSquareGrid{} samples on a side, \RcOneSquare{}~mm, z up; raw scan, no ink detector. It is the square of Figure~\ref{fig:r2c}~(a), with crossings between traced surfaces adjudicated. The surface is interpolated bilinearly between the nodes of the tracer's own grid and the scan trilinearly between its voxels; \FullSquareMissing{} samples fall in a chunk missing from the scan. The grey levels are stretched linearly between the percentiles \FullSquareLo{} and \FullSquareHi{} of the samples. As in Figure~\ref{fig:r2c}~(a), the tracer's grid lies about \RcOneRotation{} degrees in plane from the nearest axis aligned grid, so the fibres run slanted. No coordinate on the scroll is given. The scale bar is \FullSquareBar{}~mm. Drawn by \texttt{src/tools/c-f25-square-full.py}.}
\newcommand{\CaptionVeinW}{Youssef Nader's v8-in ink model on the labelled region w016 of PHerc. 0139: (a) on the organisers' surface, (b) on our sheet Bx, each in its own grid and in the depth order measured on w016, probability from black (none) to white (certain). The sky blue outline is the labelled ink among the clean held out label points matched to each surface within 4 voxels. On those points the balanced accuracy is 0.6305 on the organisers' surface (171,693 points) and 0.6017 on our sheet (91,075 points). These are the maps of the calibration read on every matched point; the depth order of Section~\ref{sec:checks} was measured on the same points of another render of that surface, read at a coarser stride, so its value for the organisers' surface differs. Panel (b) is turned by a quarter turn. The scale bars are 1 mm. Drawn by \texttt{src/tools/c-f22-v8in-w016.py}.}
\newcommand{\CaptionVeinSquare}{Youssef Nader's v8-in ink model on the headline surface of PHerc.~0826, grown from the centre of the certified square of seed 6273. It covers that surface's square as it stood before crossings between traced surfaces were adjudicated, 29.8961 mm (panels b and c of Figure~\ref{fig:r2c}), on which the shares of Section~\ref{sec:checks} are counted. It is drawn in the square's own grid, turned as in Figure~\ref{fig:r2c}, in both depth orders: the top row in the order measured on w016, the bottom row in the other order. (a) and (d) the sheet; (b), (c), (e) and (f) its two null copies, the same surface moved along its normal into the dark gap on either side and read the same way. Probability from black (none) to white (certain), one stretch for every panel. We, the authors, looked at the six maps by eye and see no letter forms on the sheet. The scale bar is 5 mm. No coordinate on the scroll is given. Drawn by \texttt{src/tools/c-f23-v8in-0826.py}.}
\newcommand{\CaptionVeinMore}{The other reads of Youssef Nader's v8-in ink model on PHerc.~0826, sheet only, each panel in its own grid, in the depth order measured on w016 and in the other order: (a) and (b) the square grown from the start of seed 5364, as it stood before crossings between traced surfaces were adjudicated; (c) and (d) the whole surface that the organisers' tracer grew on route R2d from sheet 0 of seed 2715, whose null copies were not read. Each panel is labelled with its seed and depth order. Probability from black (none) to white (certain), one stretch for every panel; white in (c) and (d) is where no tile of the read reached. Part of the seed 2715 surface runs through air (Section~\ref{sec:checks}), and there the map is noise. We, the authors, looked at the four maps by eye and see no letter forms in any of them. The numbers of these reads are those of Table~\ref{tab:veinreads}. The scale bars are 5 mm. No coordinate on the scroll is given. Drawn by \texttt{src/tools/c-f24-v8in-more.py}.}
